(1) 定常方程 \(\mu u''=0\)，边界为 \(u(0)=b u'(0)\)、\(U-u(H)=b u'(H)\)。解 \(u(y)=U(y+b)/(H+2b)\)。流体实际接触速度为 \(u(0)=Ub/(H+2b)\) 与 \(u(H)=U(H+b)/(H+2b)\)，两侧相对壁面滑移幅度相同。

(2) 剪应力大小 \(\tau_w=\mu U/(H+2b)\)。体内耗散每面积 \(\mathcal D_{\rm bulk}=\int_0^H\mu(u')^2dy=\mu U^2H/(H+2b)^2\)。界面滑移耗散为 \(2\tau_w^2b/\mu=2\mu U^2b/(H+2b)^2\)。两者相加为 \(\tau_w U\)，与上壁输入功率严格相等。

(3) 由 \(U=\tau_w(H+2b)/\mu\) 得内部速度差 \(u(H)-u(0)=\tau_w H/\mu\)，双壁滑移总量 \(2b\tau_w/\mu\)。当 \(b\to0\) 是经典无滑移Couette；在固定 \(\tau_w\) 下 \(b\to\infty\) 所需板速发散，表明该极限不能同时假定固定板速和固定外力。

(4) 对于 \(G=-p_x\) 有 \(\mu u''=-G\)。附加满足齐次滑移边界的流速 \(u_G(y)=G[y(H-y)+bH]/(2\mu)\)；总 \(u=u_{\rm Couette}+u_G\)，每单位宽度流量 \(Q=UH/2+GH^3/(12\mu)+GbH^2/(2\mu)\)。附加 \(u_G'(0)=GH/(2\mu)\)，\(u_G'(H)=-GH/(2\mu)\)，满足两壁各自正确的法向滑移符号。
